% Section 2, part 1 of 8; label prefix: set. The main file inputs only this part; each part inputs the next.
% Part 1: introduction; the formal quintic, genus-zero data, PF5 = z^5 D PF4 - 1, the quoted inputs (M) and (G0).
% Part 2: the large-radius cyclic structure (statement; the proof of (i)-(ii) is in the appendix on the large-radius cyclic structure).
% Part 3: its corollary; points; Frobenius data.
% Part 4: the R-matrix and its rows, (QRR), the twist sign, the rotation of h_w; the graph sum, (T).
% Part 5: the Frobenius basis at the conifold point, the conifold frames, the conifold expansion, tau, the canonical coordinates.
% Part 6: the model, its rows and the T-system, M_0 and its symplecticity.
% Part 7: the model cohomological field theory, its homogeneity, the dilaton lemma.
% Part 8: the hatted theory, the outer graph sum; Property N(g) and the constants c_g.
% No new macros. Local letters are declared where they occur.
\section{The formal quintic, the model and the outer graph sum}
\label{sec:set}

This section fixes the objects of the note and states the quoted inputs (M), (G0) (\cref{sec:set-fq}), (QRR) (\cref{sec:set-rows}) and (T) (\cref{sec:set-graph}). Four facts used in the definitions are proved later, without using the definitions that cite them: \cref{lem:gau-confexp}, \cref{cor:gau-rowsNN}, \cref{lem:gau-colrows}, and the symplecticity of $R'$ (from \cref{thm:gam-D}, whose proof does not use $R'$). Printed numbers are values of the definitions; no finite computation enters a proof.

%-------------------------------------------------------------------------------
\subsection{The formal quintic and its genus-zero data}
\label{sec:set-fq}

Throughout, $\zeta\coloneqq\re^{2\pi\ri/5}$, colours $i,j$ and flat indices $a,b$ run over $\bbZ/5$, and $\overline{\mathcal M}_{g,n}$ is the moduli space of stable curves of genus $g$ with $n$ marked points.

\begin{defn}[The formal quintic]
\label{def:set-fq}
Let $\bbT=(\bbC^*)^5$ act on $\bbP^4$ with weights $\lambda_0,\dots,\lambda_4$ on $H^0(\bbP^4,\cO(1))$, with fixed points $p_0,\dots,p_4$, and let $H\coloneqq c_1^{\bbT}(\cO(1))$; so $H|_{p_i}=\lambda_i$, the tangent weights at $p_i$ are $\lambda_i-\lambda_j$ ($j\neq i$), and $\cO(5)|_{p_i}$ has the weight $5\lambda_i$ \cite[(10)]{LP20}. For $g,d\geq0$ the \emph{formal quintic invariant} is $N_{g,d}\coloneqq\int_{[\overline{\mathcal M}_g(\bbP^4,d)]^{\mathrm{vir}}}e\bigl(R\pi_*f^*\cO(5)\bigr)$ \cite[Section~0.2, (5)]{LP18}, \cite[Section~0.2, (5)]{LP20} (whenever the moduli space is non-empty), defined by virtual localisation \cite{MR1666787}, the Euler class being taken on each fixed locus as the quotient of the top Chern classes of the two terms of $R\pi_*f^*\cO(5)$. It is a rational function of the weights, homogeneous of degree $0$ (the rank $5d+1-g$ equals the virtual dimension), specialised at $\lambda_i=\zeta^i\lambda$, where it does not depend on $\lambda$; we put $\lambda=1$. (The second term, $\bbE^\vee\otimes\cO(5)|_{p_i}$ on the contracted components, has an invertible top Chern class there, as $5\zeta^i\neq0$. The other denominators of the localisation formula are products of the weights $(a\lambda_i+b\lambda_j)/d-\lambda_k$ ($a+b=d$, non-zero before the specialisation), $(\lambda_i-\lambda_j)/d_1+(\lambda_i-\lambda_k)/d_2$ and of classes $(\lambda_i-\lambda_j)/d-\psi$ ($j,k\neq i$); none vanishes at $\lambda_i=\zeta^i$, as a proper convex combination of two distinct fifth roots of unity has modulus $<1$.) This is the specialisation of \cite[Section~0.2]{LP18} and \cite[Section~0.2]{LP20}; see \cite[Section~3.1; Section~3.2, (9)]{GJR18}. The resulting hypergeometric series (\cref{lem:set-periods}) is studied in \cite{ZZ08}. For $g\geq2$ the \emph{genus-$g$ potential} is $F_g\coloneqq\sum_{d\geq0}N_{g,d}\,\Qn^d\in\bbQ\fps{\Qn}$, the term $d=0$ included; along the mirror map $\Qn=\re^{T(q)}\in q\bbQ\fps q$ (\cref{def:set-genus0}) it becomes a power series in $q$, again written $F_g$.

More generally, for $2g-2+n>0$ and $\gamma_1,\dots,\gamma_n\in H^*_{\bbT}(\bbP^4)$, with the stabilisation map $\rho$,
\[
\Omega_{g,n}(\gamma_1,\dots,\gamma_n)\coloneqq\sum_{d\geq0}\Qn^d\,\rho_*\Bigl(e\bigl(R\pi_*f^*\cO(5)\bigr)\prod_{j}\mathrm{ev}_j^*\gamma_j\cap[\overline{\mathcal M}_{g,n}(\bbP^4,d)]^{\mathrm{vir}}\Bigr)\in H^*(\overline{\mathcal M}_{g,n})\fps{\Qn},
\]
defined and specialised in the same way; with the \emph{twisted pairing} $\eta(a,b)\coloneqq\int_{\bbP^4}ab\,e_{\bbT}(\cO(5))$ these are the classes of the cohomological field theory of the formal quintic \cite[Section~4.1]{GJR18}, and $F_g=\int_{\overline{\mathcal M}_g}\Omega_{g,0}$. At $\lambda_i=\zeta^i$, $\prod_j(H-\lambda_j)=0$ becomes $H^5=1$, and the \emph{state space} is $\bbH\coloneqq\bbQ[H]/(H^5-1)$, with the flat basis $H^a$ ($a\in\bbZ/5$), $H^0=1$. Restriction to $p_i$ is $H\mapsto\zeta^i$, and an element of $\bbH\otimes\bbQ(\zeta)$ is determined by its five restrictions. As $e_{\bbT}(\cO(5))=5H$ and $e_{\bbT}(T_{p_i}\bbP^4)=\prod_{j\neq i}(\lambda_i-\lambda_j)=5\zeta^{4i}$, localisation gives
\[
\eta(H^a,H^b)=\sum_{i\in\bbZ/5}\frac{5\,\zeta^{i(a+b+1)}}{5\,\zeta^{4i}}=5\,\delta_{a+b\equiv3\bmod5},\qquad\eta^{cc'}=\tfrac15\,\delta_{c+c'\equiv3\bmod5},
\]
$(\eta^{cc'})$ the inverse matrix. The $N_{g,d}$ are rational: permuting the coordinates of $\bbC^5$ permutes the weights and preserves $H$ and $\cO(5)$ with their linearisations, so the rational function is symmetric, and its value at $\lambda_i=\zeta^i$ is fixed by the Galois group of $\bbQ(\zeta)$.
\end{defn}

\begin{defn}[Genus-zero data]
\label{def:set-genus0}
Let $\Dq\coloneqq q\,\dd/\dd q$, and let $\logq$ be an indeterminate with $\Dq\logq=1$, a formal $\log q$; $\Dq$ acts on $\bbQ\fps q[\logq]$, and $z$ is a constant for $\Dq$. Put
\[
C_0\coloneqq\sum_{d\geq0}\frac{(5d)!}{(d!)^5}q^d,\qquad\tilde I_1\coloneqq5\sum_{d\geq1}\frac{(5d)!}{(d!)^5}\Bigl(\sum_{m=d+1}^{5d}\frac1m\Bigr)q^d,\qquad I_1\coloneqq C_0\logq+\tilde I_1,
\]
the \emph{mirror map} $T\coloneqq I_1/C_0=\logq+\tilde I_1/C_0$, $C_1\coloneqq\Dq T$, and $\re^T\coloneqq q\exp(\tilde I_1/C_0)\in q\bbQ\fps q$; thus $C_0,C_1\in1+q\bbQ\fps q$. Put $L\coloneqq(1-5^5q)^{-1/5}$, so that $\Dq L=5^4qL^6=(L^6-L)/5$. The \emph{I-function} $I(q,z)$ is the $\bbH\otimes\bbQ(\zeta)$-valued series, $z$ times a series in $z^{-1}$ with coefficients in $\bbQ(\zeta)\fps q[\logq]$, with the restrictions
\[
I|_{p_i}=I_i(q,z)\coloneqq z\,q^{\zeta^i/z}\sum_{d\geq0}q^d\,\frac{\prod_{k=1}^{5d}(5\zeta^i+kz)}{\prod_{j\in\bbZ/5}\prod_{k=1}^{d}(\zeta^i-\zeta^j+kz)},\qquad q^{\zeta^i/z}\coloneqq\exp(\zeta^i\logq/z).
\]
The fifth-order and the fourth-order operators are
\begin{equation}
\label{eq:set-PF5}
\PFfive\coloneqq(z\Dq)^5-1-q\prod_{k=1}^5(5z\Dq+kz),
\end{equation}
\begin{equation}
\label{eq:set-PF4}
\PFfour\coloneqq\Dq^4-5q\prod_{k=1}^4(5\Dq+k)=(1-5^5q)\Dq^4-6250q\Dq^3-4375q\Dq^2-1250q\Dq-120q .
\end{equation}
With $u\coloneqq1-5^5q$ (\cref{def:set-points}), $\PFfour=u\Dq^4-(1-u)\bigl(2\Dq^3+\frac75\Dq^2+\frac25\Dq+\frac{24}{625}\bigr)$, the coefficients being the elementary symmetric functions of $\frac15,\dots,\frac45$. The eigenvalues of $H*$, the canonical coordinates and the normalisations $\Delta_i$ are consequences of (M) and (G0); they are stated in \cref{cor:set-cyclicPF5}~(a).
\end{defn}

\begin{lem}[$\PFfive=z^5\Dq\PFfour-1$]\genusstatus{AI-proposed, unchecked}
\label{lem:set-pf5}
As operators on any differential $\bbQ$-algebra that contains $q$ with $\Dq q=q$ and in which $z$ is a constant (for instance any differential $\bbQ\fps q$-algebra), $\PFfive=z^5\,\Dq\circ\PFfour-1$.
\end{lem}

\begin{proof}
As $\Dq\circ q=q\circ(\Dq+1)$, $\Dq\circ\PFfour=\Dq^5-q\prod_{k=1}^5(5\Dq+k)$; and as $z$ is a constant, $\PFfive=z^5\bigl(\Dq^5-q\prod_{k=1}^5(5\Dq+k)\bigr)-1$.
\end{proof}

\begin{thm}[The genus-zero mirror theorem (M); quoted]\genusstatus{AI-proposed, unchecked}
\label{thm:set-mirror}
For $d\geq0$ and $n\geq1$, with $n\geq3$ if $d=0$, let $E_{0,n,d}\coloneqq R^0\pi_*f^*\cO(5)$ on $\overline{\mathcal M}_{0,n}(\bbP^4,d)$, a vector bundle of rank $5d+1$ because $R^1\pi_*f^*\cO(5)=0$ in genus $0$, and let
\[
\langle\psi^{k_1}\gamma_1,\dots,\psi^{k_n}\gamma_n\rangle_{0,n,d}\coloneqq\int_{[\overline{\mathcal M}_{0,n}(\bbP^4,d)]}e_{\bbT}(E_{0,n,d})\prod_{j=1}^n\mathrm{ev}_j^*\gamma_j\,\psi_j^{k_j}
\]
be the genus-zero twisted invariants. The \emph{J-function}, in the normalisation of \cite[Section~6, (10)]{MR2276766} with the twisted pairing, is
\[
J(t,z)\coloneqq z+t+\sum_{n,d}\frac{\Qn^d}{n!}\sum_{a\in\bbZ/5}\Bigl\langle t,\dots,t,\frac{H^a}{z-\psi}\Bigr\rangle_{0,n+1,d}H^\vee_a,\qquad t\in H^*_{\bbT}(\bbP^4),
\]
where $H^\vee_a$ is the $\eta$-dual basis and the sum runs over the $(n,d)$ for which $\overline{\mathcal M}_{0,n+1}(\bbP^4,d)$ is defined. For $t=\tau H$ the divisor equation (\cref{thm:set-G0}~(b)) gives $J(\tau H,z)=\re^{\tau H/z}\bigl(z+\sum_{d\geq1}(\Qn\re^{\tau})^d\sum_a\langle H^a/(z-\psi)\rangle_{0,1,d}H^\vee_a\bigr)$. Since $T=\logq+\tilde I_1/C_0$ contains the formal logarithm, $J(T(q)H,z)$ at $\Qn=1$ is \emph{defined} by this formula, with $\re^{dT}\coloneqq(\re^T)^d$, $\re^T=q\exp(\tilde I_1/C_0)$ (\cref{def:set-genus0}), and $\re^{TH/z}\coloneqq\sum_{k\geq0}(TH)^k/(k!\,z^k)$:
\[
J(TH,z)\coloneqq\re^{TH/z}\Bigl(z+\sum_{d\geq1}\re^{dT}\sum_{a\in\bbZ/5}\Bigl\langle\frac{H^a}{z-\psi}\Bigr\rangle_{0,1,d}H^\vee_a\Bigr).
\]
Then at $\lambda_i=\zeta^i$ (see also (d) of the translation below),
\[
\text{(M)}\qquad J\bigl(T(q)H,z\bigr)=\frac{I(q,z)}{C_0(q)},
\]
an identity of $z$ times series in $z^{-1}$ with coefficients in $\bbH\otimes\bbQ(\zeta)\fps q[\logq]$.
\end{thm}

\emph{The pin.} (M) is \cite[Theorem~11.6]{MR1408320} (set-up in Section~9, pp.~23--27; the proof rests on Proposition~11.4) in the Calabi--Yau case $n=4$, $r=1$, $l_1=5$, $\lambda'=0$, at $\lambda_i=\zeta^i$. In this paragraph the letters are Givental's: $\bbC\bbP^n$ with the torus of $\bbC^{n+1}$ (weights $\lambda_\alpha$), the bundle $\oplus_a\cO(l_a)$ with fibre weights $\lambda'_a$, the space $Y_{2,d}$ of genus-zero stable maps of degree $d$ with markings $x_0,x_1$, the Euler class $E'_d$ of the bundle of sections vanishing at $x_0$, $c=-\psi_0$, and for $r=1$ (p.~27) $Z_i\coloneqq\sum_{d\geq0}q^d\int_{Y_{2,d}}e_0^*\bigl(\prod_{j\neq i}(p-\lambda_j)\bigr)(\hbar+c)^{-1}E'_d$, the term $d=0$ being $1$. With $\mathfrak f\coloneqq\sum_d\prod_a(l_ad)!/(d!)^{n+1}q^d$ and $\mathfrak g_l\coloneqq\sum_{d\geq1}q^d\prod_a(l_ad)!/(d!)^{n+1}\sum_{m=1}^{ld}\frac1m$, Theorem~11.6 states, for $\sum_al_a=n+1$, that
\[
Z_i^*\coloneqq\sum_{d\geq0}q^d\,\frac{\prod_a\prod_{m=1}^{l_ad}(l_a\lambda_i-\lambda'_a+m\hbar)}{\prod_\alpha\prod_{m=1}^d(\lambda_i-\lambda_\alpha+m\hbar)}
\]
coincides with $Z_i$ after (1) the substitution $Q=q\exp\{\sum_al_a(\mathfrak g_{l_a}-\mathfrak g_1)/\mathfrak f\}$, (2) the multiplication of $Z_i(Q(q),\hbar)$ by $\mathcal E_i\coloneqq\exp\{(\mathfrak f\hbar)^{-1}[\sum_a(l_a\lambda_i-\lambda'_a)\mathfrak g_{l_a}-(\sum_\alpha(\lambda_i-\lambda_\alpha))\mathfrak g_1]\}$, and (3) the multiplication by $\mathfrak f$. With $\mathfrak s_i^*(t)\coloneqq\re^{\lambda_it/\hbar}Z_i^*(\re^t,\hbar)$, $\mathfrak s_i(T)\coloneqq\re^{\lambda_iT/\hbar}Z_i(\re^T,\hbar)$ and $\re^T=Q(\re^t)$, this reads $\mathfrak s_i^*(t)=\mathfrak f\,\mathcal E_i\,\re^{-\lambda_i(T-t)/\hbar}\,\mathfrak s_i(T)$; as $T-t=\sum_al_a(\mathfrak g_{l_a}-\mathfrak g_1)/\mathfrak f$ and $\sum_al_a=n+1$, the terms in $\lambda_i$ cancel in the exponent, and $\mathfrak s_i^*(t)=\mathfrak f\exp\{(\mathfrak g_1\sum_\alpha\lambda_\alpha-\sum_a\lambda'_a\mathfrak g_{l_a})/(\hbar\mathfrak f)\}\,\mathfrak s_i(T)$. (\cite[Corollary~11.7]{MR1408320} prints $l_\alpha$ for $\lambda_\alpha$ here.)

\emph{Translation.} Put $n=4$, $r=1$, $l_1=5$, $\lambda'_1=0$, $p=H$, $\hbar=z$, $\re^t=q$, $t=\logq$. (a) $\mathfrak f=C_0$ and $5(\mathfrak g_5-\mathfrak g_1)=\tilde I_1$, so (1) is the mirror map of \cref{def:set-genus0}. (b) The factor above is $C_0\exp\{e_1(\lambda)\mathfrak g_1/(zC_0)\}$, $e_1(\lambda)=\sum_\alpha\lambda_\alpha$, which is $C_0$ at $\lambda_i=\zeta^i$. (c) $\mathfrak s_i^*(t)=I_i(q,z)/z$. (d) $\mathfrak s_i(T)=J(TH,z)|_{p_i}/z$: the $\eta$-dual of $\gamma\mapsto\int_{\overline{\mathcal M}_{0,1}(\bbP^4,d)}\beta\,\mathrm{ev}^*\gamma$ is $\mathrm{ev}_*(\beta)/e_{\bbT}(\cO(5))$, and $E_{0,1,d}$ is an extension of $\mathrm{ev}^*\cO(5)$ by the sections vanishing at the marking ($f^*\cO(5)$ is globally generated), so the definition of $J(TH,z)$ reads $J(TH,z)=\re^{TH/z}\bigl(z+\sum_{d\geq1}\re^{dT}\mathrm{ev}_*\bigl(E'_d/(z-\psi)\bigr)\bigr)$; on the other side $E'_d$ is pulled back along the map forgetting $x_1$, the string equation gives $\pi_*(\hbar-\psi_0)^{-1}=\hbar^{-1}(\hbar-\psi_0)^{-1}$ ($d\geq1$), and $\int_{\bbP^4}\prod_{j\neq i}(H-\lambda_j)\,\gamma=\gamma|_{p_i}$, so $Z_i(q,z)=1+z^{-1}\sum_{d\geq1}q^d\,\mathrm{ev}_*\bigl(E'_d/(z-\psi)\bigr)|_{p_i}$. By (a)--(d) the relation reads $C_0\,J(TH,z)|_{p_i}=I_i(q,z)$ for every $i$, which is (M).

\emph{Hypotheses.} The case is Calabi--Yau, $l_1=5=n+1$; the classes are polynomial in $\lambda$, $\lambda'$ and $\hbar$, so $\lambda'=0$ is a specialisation; and in genus zero $H^1(C,f^*\cO(5))=0$, so Givental's class $E_d$ is the Euler class of $E_{0,n,d}$ (convexity is all that is needed). After removing $\re^{\lambda_iT/z}$ both sides are, coefficientwise in $q$ and $z^{-1}$, polynomials in the weights, while $T$ and $C_0$ do not depend on them; so the specialisation $\lambda_i=\zeta^i$ is a ring homomorphism, and $e_1(\zeta^0,\dots,\zeta^4)=0$ gives (M).

\begin{thm}[The genus-zero axioms (G0); quoted]\genusstatus{AI-proposed, unchecked}
\label{thm:set-G0}
At $\lambda_i=\zeta^i$ the classes $\Omega_{0,n}$ of \cref{def:set-fq} form the genus-zero part of a cohomological field theory with unit $1$ for the pairing $\eta$ \cite[Section~4.1]{GJR18}, \cite[Def.~0.3]{PPZ15}, and the following hold.
\begin{itemize}
\item[(a)] For $d\geq1$ the invariant $\langle H^{a_1},\dots,H^{a_n}\rangle_{0,n,d}$, before the specialisation, lies in $\bbQ[\lambda_0,\dots,\lambda_4]$, is symmetric in the $\lambda_j$ and homogeneous of degree $\sum_ja_j-n$: the moduli space is proper of dimension $5d+n+1$, $E_{0,n,d}$ has rank $5d+1$, and a permutation of the coordinates of $\bbC^5$ permutes the weights and preserves $H$ and $\cO(5)$ with their linearisations.
\item[(b)] The string and divisor equations hold, also with descendant insertions at the other points (with the usual correction terms, from $\psi_j=\pi^*\psi_j+D_j$): $E_{0,n+1,d}$ is the pull-back of $E_{0,n,d}$ along the map $\pi$ forgetting a point, and $\int_dH=d$.
\item[(c)] The splitting axiom for $\eta$ holds (by the normalisation sequence along a boundary divisor, $e(E)$ restricted to the divisor times $e(\mathrm{ev}_{\mathrm{node}}^*\cO(5))$ is $e(E_1)e(E_2)$), hence the WDVV equations, and so does the genus-zero topological recursion relation.
\item[(d)] The quantum differential equation \cite[Section~5.2, the display before (14), of which (14) is the form in the $q$-direction]{GJR18}: let
\[
\langle\!\langle x,y,w\rangle\!\rangle(t)\coloneqq\sum_{n,d}\frac{\Qn^d}{n!}\langle x,y,w,t,\dots,t\rangle_{0,n+3,d},\qquad x*y\coloneqq\sum_{c,c'}\langle\!\langle x,y,H^c\rangle\!\rangle\,\eta^{cc'}H^{c'},
\]
the quantum product at $t$, $(\eta^{cc'})$ being the inverse matrix of $(\eta(H^c,H^{c'}))$, and let $J_a\coloneqq\partial J/\partial t^a$ be the derivative in the direction $H^a$. Then $z\,\partial_{t^b}J_a=\sum_c(H^b*H^a)^cJ_c$, $J_a=H^a+O(z^{-1})$, and $J_0=J/z$ (string equation). At $t=T(q)H$, $\Qn=1$ these series are understood as in \cref{thm:set-mirror}: by (b), $\langle\!\langle x,y,w\rangle\!\rangle(TH)\coloneqq\eta(xy,w)+\sum_{d\geq1}\re^{dT}\langle x,y,w\rangle_{0,3,d}$, and $J_a(TH,z)$ is $\partial J/\partial t^a$ at $t=\tau H$ with $(\Qn\re^{\tau})^d$ replaced by $\re^{dT}$ and $\re^{\tau H/z}$ by $\re^{TH/z}$. (d) follows from (b), (c) and the genus-zero topological recursion relation.
\end{itemize}
\end{thm}

% Section 2, part 2 of 8: the large-radius cyclic structure (statement).

Along the line $t=T(q)H$, with $\Qn=1$, functions of $t$ are functions of $q$: by the divisor equation (G0)(b) the insertions of $TH$ are resummed, so that the functions lie in $\bbQ[T]\fps{\re^T}$, and $T\mapsto T(q)$, $\re^T\mapsto q\exp(\tilde I_1/C_0)$ (\cref{def:set-genus0}) maps them to $\bbQ\fps q[\logq]$, intertwining $\partial_T$ with $C_1^{-1}\Dq$, and $\partial_T\coloneqq C_1^{-1}\Dq$ is the derivative in the coordinate $T$ of $H$, i.e.\ $\partial_{t^1}$ along the line (not the derivation $\de_t$ of \cref{def:gau-conc}). By (G0), the quantum product $*$ at $t=TH$ is commutative and associative, has the unit $1$, and is Frobenius for $\eta$: $\eta(x*y,w)=\langle\!\langle x,y,w\rangle\!\rangle=\eta(x,y*w)$.

\begin{prop}[The large-radius cyclic structure]\genusstatus{AI-proposed, unchecked}
\label{prop:set-cyclicLR}
Assume (M) and (G0), and for (iii) the flatness part of (QRR), \cref{def:set-rows}~(b). Along the line $t=T(q)H$, at $\lambda_i=\zeta^i$ and $\Qn=1$, the following hold.
\begin{enumerate}
\item In the flat basis, $H*H^a=\kappa^{\mathrm{LR}}_aH^{a+1}$ for $a\in\bbZ/5$, where
\[
\kappa^{\mathrm{LR}}=(\kappa^{\mathrm{LR}}_0,\dots,\kappa^{\mathrm{LR}}_4)=\Bigl(1,\ \frac{L^5}{C_0^2C_1^3},\ 1,\ \frac{C_0}{C_1},\ \frac{C_0}{C_1}\Bigr).
\]
\item Consequently, for every solution $v$ of the quantum differential equation $z\partial_Tv=H*v$, in any differential $\bbQ\fps q$-algebra in which $z$ is a constant, the lowered components $v_a\coloneqq\eta(H^a,v)$ satisfy $z\Dq v_a=C_1\kappa^{\mathrm{LR}}_a\,v_{a+1}$ for all $a\in\bbZ/5$.
\item For every colour $i$ the rows of \cref{def:set-rows} (defined in \cref{sec:set-rows}) satisfy $(z\Dq+\zeta^iL)\Qr^{(i)}_a=C_1\kappa^{\mathrm{LR}}_a\,\Qr^{(i)}_{a+1}$ for $a\in\bbZ/5$. Equivalently, for $X_a\coloneqq\zeta^{-i}\Qr^{(i)}_a$,
\[
X_0=L\Phi_i/C_0,\qquad X_{a+1}=(C_1\kappa^{\mathrm{LR}}_a)^{-1}(z\Dq+\zeta^iL)X_a\quad(a=0,\dots,3),\qquad(z\Dq+\zeta^iL)X_4=C_0X_0 .
\]
\end{enumerate}
\end{prop}

The superscript LR distinguishes $\kappa^{\mathrm{LR}}$ from the coefficients $\kappa_a$ of \cref{prop:gau-quintic}. Parts (i) and (ii) are proved in \cref{sec:set-lrproof}, (iii) after \cref{def:set-rows}. Below, $\Mcyc^{\mathrm{LR}}\coloneqq\partial_T^2\circ(\kappa^{\mathrm{LR}}_1)^{-1}\circ\partial_T^2\circ C_0^{-1}$, and $\Mcyc^{\mathrm{LR}}=\kappa^{\mathrm{LR}}_3\,\PFfour$ by \eqref{eq:set-McycLR}.

% Section 2, part 3 of 8: the corollary of prop:set-cyclicLR; points; Frobenius data.

\begin{cor}[Eigenvalues, idempotents and the fifth-order equation]\genusstatus{AI-proposed, unchecked}
\label{cor:set-cyclicPF5}
Assume (M) and (G0), and for the last sentence of (b) the flatness part of (QRR), \cref{def:set-rows}~(b).
\begin{itemize}
\item[(a)] $(H*)^5=(L/C_1)^5$. The eigenvalues of $H*$ are $\mu_i\coloneqq\zeta^iL/C_1$ ($i\in\bbZ/5$), labelled by $\mu_i(0)=\zeta^i$, the classical eigenvalue at $p_i$. The idempotents are
\[
e_i=\frac15\sum_{a=0}^4\Bigl(\prod_{b<a}\kappa^{\mathrm{LR}}_b\Bigr)\mu_i^{-a}H^a,\qquad\Delta_i\coloneqq\frac1{\eta(e_i,e_i)}=\frac{C_0^2\,\zeta^{3i}}{L^2}.
\]
The canonical coordinates along the line are functions $u_i$ with $\partial_Tu_i=\mu_i$, that is $\Dq u_i=\zeta^iL$; they are determined up to additive constants, the constants of $\Dq$ on $\bbQ(\zeta)\fps q[\logq]$ being $\bbQ(\zeta)$ (as in the proof of \cref{lem:set-periods}; for the conifold chart see \cref{lem:set-canon}).
\item[(b)] For every solution $v$ of $z\partial_Tv=H*v$ in a differential $\bbQ\fps q$-algebra in which $z$ is a constant, $C_0v_0$ solves $\PFfive$. In particular, for the column $v=R(z)\tilde e_i\re^{u_i/z}$ of \cref{def:set-rows}, the function $\re^{u_i/z}\,z\,L\,\Phi_i(\zh,u)$ solves $\PFfive$: this is the WKB form of colour $i$.
\end{itemize}
\end{cor}

\begin{proof}
(a) By \cref{prop:set-cyclicLR}~(i), $(H*)^5H^a=\prod_b\kappa^{\mathrm{LR}}_b\,H^a=(L/C_1)^5H^a$, so the characteristic polynomial of $H*$ is $x^5-(L/C_1)^5$, with the pairwise distinct roots $\zeta^iL/C_1$, as $L/C_1\in1+q\bbQ\fps q$. The equation $H*x=\mu_ix$ for $x=\sum_ax_aH^a$ reads $x_a=x_{a-1}\kappa^{\mathrm{LR}}_{a-1}/\mu_i$, consistently around the cycle as $\prod_b\kappa^{\mathrm{LR}}_b=\mu_i^5$. As $*$ is commutative and associative with unit $1$, the eigenlines are ideals spanned by idempotents $e_i$ with $\sum_ie_i=1$; for the coefficient $x_0^{(i)}$ of $H^0$ in $e_i$, the coefficient of $H^a$ in $\sum_ie_i=1$ is $\bigl(\prod_{b<a}\kappa^{\mathrm{LR}}_b\bigr)(C_1/L)^a\sum_ix_0^{(i)}\zeta^{-ia}=\delta_{a0}$, whose unique solution is $x_0^{(i)}=1/5$. By the Frobenius property and $\eta(1,H^a)=5\,\delta_{a3}$, $\eta(e_i,e_i)=\eta(1,e_i)=\kappa^{\mathrm{LR}}_0\kappa^{\mathrm{LR}}_1\kappa^{\mathrm{LR}}_2\,\mu_i^{-3}=L^2/(C_0^2\zeta^{3i})$. Finally $\partial_Tu_i=\mu_i$ and $\Dq=C_1\partial_T$ give $\Dq u_i=\zeta^iL$.

(b) By \cref{prop:set-cyclicLR}~(ii), $v_1=z\partial_Tv_0$, $v_2=z^2(\kappa^{\mathrm{LR}}_1)^{-1}\partial_T^2v_0$, $v_3=z^3\partial_T\bigl((\kappa^{\mathrm{LR}}_1)^{-1}\partial_T^2v_0\bigr)$ and, by \eqref{eq:set-McycLR},
\[
v_4=z^4(\kappa^{\mathrm{LR}}_3)^{-1}\Mcyc^{\mathrm{LR}}(C_0v_0)=z^4\,\PFfour(C_0v_0).
\]
The closing relation $z\partial_Tv_4=\kappa^{\mathrm{LR}}_4v_0$ becomes $z^5C_1^{-1}\Dq\PFfour(C_0v_0)=(C_0/C_1)v_0$, i.e.\ $\PFfive(C_0v_0)=0$ by \cref{lem:set-pf5}. The last statement concerns the rows and is proved after \cref{def:set-rows}.
\end{proof}

\begin{defn}[Points and variables]
\label{def:set-points}
Put $u\coloneqq1-5^5q$ and $s\coloneqq u^{1/5}$, so that $L=u^{-1/5}=s^{-1}$ and $\Dq=q\,\dd/\dd q=-(1-u)\,\dd/\dd u$, and put $\zh\coloneqq z/s^4=L^4z$, the hatted variable. The \emph{large-radius point} is $q=0$ ($u=1$, $L=s=1$, $\zh=z$), the \emph{conifold point} is $u=0$ ($q=5^{-5}$), and the \emph{orbifold point} is $q=\infty$ ($L=0$). Near $q=0$ the functions of the note lie in $\bbQ\fps q[\logq]$; near the conifold point the note uses their conifold expansions (\cref{def:set-confexp}), which are Laurent series in $s$, with $\Dq=-\frac{1-s^5}5s^{-4}\frac{\dd}{\dd s}$. The conifold expansion substitutes the periods of a conifold frame for $C_0$ and $I_1$ in a representative in the ring $\Nring$; it is a formal change of frame, and no analytic continuation is involved.
\end{defn}

\begin{defn}[Frobenius data and the topological field theory]
\label{def:set-frob}
Along the line $t=T(q)H$ let $e_i$ be the idempotents and $\Delta_i=1/\eta(e_i,e_i)$ (\cref{cor:set-cyclicPF5}~(a)), $\eta$ the pairing of \cref{def:set-fq}. The \emph{normalised canonical basis} is $\tilde e_i\coloneqq\Delta_i^{1/2}e_i$, so $\eta(\tilde e_i,\tilde e_j)=\delta_{ij}$, with the square roots $\zeta^{1/2}\coloneqq\zeta^3$ and $\Delta_i^{1/2}\coloneqq\zeta^{4i}\,C_0/L\in\bbQ(\zeta)\fps q$, so that $(\Delta_i^{1/2})^2=\zeta^{8i}C_0^2/L^2=\Delta_i$. At $q=0$, $\Delta_i(0)=W^0_i\coloneqq\prod_{j\neq i}(\lambda_i-\lambda_j)/(5\lambda_i)=\zeta^{3i}$. The \emph{topological field theory of the formal quintic} is the TFT $\TFT$ on $\bbH\otimes\bbQ(\zeta)\fps q$ with $\TFT_{g,n}(\tilde e_i,\dots,\tilde e_i)\coloneqq\Delta_i^{(2g-2+n)/2}=(\zeta^{4i}C_0/L)^{2g-2+n}$ and $\TFT_{g,n}$ vanishing on mixed colours; its unit is $1_{\TFT}\coloneqq\sum_i\Delta_i^{-1/2}\tilde e_i=\sum_ie_i=1$. We also regard $\TFT$ as a TFT on the colour space, with the orthonormal basis $\tilde e_i$ ($i\in\bbZ/5$) and the unit $1_{\TFT}$. By \cref{thm:set-GT}, since $R|_{z=0}=1$, $\TFT$ is the degree-zero part of the cohomological field theory along the line.
\end{defn}

% Section 2, part 4 of 8: the R-matrix and its rows; (QRR); the twist sign; h_w rotation; the graph sum; the two forms of the translation; (T).

%-------------------------------------------------------------------------------
\subsection{The \texorpdfstring{$R$}{R}-matrix, its rows and quantum Riemann--Roch}
\label{sec:set-rows}

\begin{defn}[The $R$-matrix and its rows]
\label{def:set-rows}
Let $R(z)=1+\sum_{k\geq1}R_kz^k$, $R_k\in\End(\bbH)\otimes\bbQ(\zeta)\fps q$, be the $R$-matrix of the semisimple cohomological field theory of the formal quintic along the line $t=T(q)H$ \cite[Section~4.1]{GJR18}, that is, the $R$-matrix of the quoted input (T), \cref{thm:set-GT}. It has the following properties: (a) is part of the quoted input (T), (b) is part of the quoted input (QRR), \cref{thm:set-qrr}, and (c) is Givental's convention:
\begin{itemize}
\item[(a)] $R$ is symplectic: $R(-z)^T\eta\,R(z)=\eta$, as matrices in the flat basis;
\item[(b)] for every colour $i$ the column $R(z)\tilde e_i\re^{u_i/z}$ solves the quantum differential equation $z\partial_Tv=H*v$, that is, $w\coloneqq R(z)\tilde e_i$ satisfies $z\partial_Tw+\mu_iw=H*w$ (\cref{cor:set-cyclicPF5}~(a));
\item[(c)] in the graph sums of \cref{def:set-graphsum}, $R$ acts with $R^{-1}(\psi)$ on the legs.
\end{itemize}
The flatness (b) and the value of $R$ at $q=0$ form the quoted input (QRR), \cref{thm:set-qrr}; that the graph sum of $R$ is the cohomological field theory of the formal quintic is the quoted input (T), \cref{thm:set-GT}, which includes (a). Property (b) is not a convention: it says that the $R$-matrix of the cohomological field theory is the asymptotic solution of the quantum differential equation, and it is used in \cref{prop:set-cyclicLR}~(iii). By \cref{lem:set-Runique}, which uses only (M) and (G0), (b) and the value at $q=0$ determine $R$, and (a) follows from them. So (QRR) and (T) together say that the graph sum of the unique series of \cref{lem:set-Runique} is the cohomological field theory of the formal quintic; no uniqueness of the $R$-matrix in (T) is used, and (a), although part of the quoted statement (T), is not an independent input.

The \emph{rows} are the lowered components $\Qr^{(i)}_a(z)\coloneqq\eta(H^a,R(z)\tilde e_i)\in\bbQ(\zeta)\fps q\fps z$, with the coefficients $\Qr^{(i)}_{a,n}\coloneqq\coef{z^n}\Qr^{(i)}_a$; $\Qr$ is the matrix $(\Qr^{(i)}_a)$, rows indexed by $a$, columns by $i$. With $\zh=L^4z$ we \emph{define} $\Phi_i(\zh,u)\coloneqq C_0\,\Qr^{(i)}_0(z)/(\zeta^iL)$, i.e.\ $\coef{\zh^n}\Phi_i\coloneqq C_0\,\Qr^{(i)}_{0,n}/(\zeta^iL^{4n+1})$, and the coefficients $\mathrm{ph}_n(u)$ by $\Phi_0(\zh,u)=\sum_{n\geq0}\mathrm{ph}_n(u)\,\zh^n$, power series in $q=(1-u)/5^5$; $\mathrm{ph}_0=1$, because $R|_{z=0}=1$ and $\eta(1,\tilde e_0)=\Delta_0^{-1/2}=L/C_0$ (\cref{def:set-frob}). The \emph{flat amplitudes} are
\[
A_n\coloneqq\coef{z^n}\bigl(C_0\,\Qr^{(0)}_0\bigr)=L^{4n+1}\,\mathrm{ph}_n(L^{-5}),\qquad\text{so that}\quad\sum_nA_nz^n=L\,\Phi_0(L^4z,L^{-5}).
\]
\end{defn}

\begin{proof}[Proof of \cref{prop:set-cyclicLR}~(iii) and of the end of \cref{cor:set-cyclicPF5}~(b)]
For $w\coloneqq R(z)\tilde e_i$, \cref{def:set-rows}~(b) and the argument of \cref{prop:set-cyclicLR}~(ii) give $z\partial_Tw_a+\mu_iw_a=\eta(H*H^a,w)=\kappa^{\mathrm{LR}}_aw_{a+1}$, $w_a=\Qr^{(i)}_a$; multiplying by $C_1$, with $C_1\mu_i=\zeta^iL$, gives $(z\Dq+\zeta^iL)\Qr^{(i)}_a=C_1\kappa^{\mathrm{LR}}_a\Qr^{(i)}_{a+1}$. Dividing by $\zeta^i$ gives the relations for $X_a$ ($C_1\kappa^{\mathrm{LR}}_4=C_0$), and $X_0=L\Phi_i/C_0$ is the definition of $\Phi_i$. For the WKB form let $\mathsf B\coloneqq\bbQ(\zeta)\fps q\lser z[\mathsf E,\mathsf E^{-1}]$ with $\Dq\mathsf E\coloneqq z^{-1}\zeta^iL\,\mathsf E$ (a local letter: $\mathsf E$ stands for $\re^{u_i/z}$, $\Dq u_i=\zeta^iL$), a differential $\bbQ\fps q$-algebra in which $z$ is a constant. By \cref{def:set-rows}~(b), $v\coloneqq R(z)\tilde e_i\,\mathsf E$ solves $z\partial_Tv=H*v$ over $\mathsf B$, as $z\partial_T(w\mathsf E)=(z\partial_Tw+\mu_iw)\mathsf E$; so $C_0v_0=\zeta^iL\Phi_i\,\mathsf E$ solves $\PFfive$ by \cref{cor:set-cyclicPF5}~(b); multiply by the constant $z\zeta^{-i}$. Moreover $\Dq(f\mathsf E)=\bigl((\Dq+z^{-1}\zeta^iL)f\bigr)\mathsf E$ gives $\PFfive(f\mathsf E)=(\PFfive^{(i)}f)\,\mathsf E$ for $f\in\bbQ(\zeta)\fps q\lser z$, where $\PFfive^{(i)}$ is $\PFfive$ with $z\Dq$ replaced by $z\Dq+\zeta^iL$; as $\mathsf E$ is a unit of $\mathsf B$, $\PFfive^{(i)}(L\Phi_i)=0$.
\end{proof}

The colour covariance $\Qr^{(i)}_{a,n}=\zeta^{i(a+1-n)}\Qr^{(0)}_{a,n}$ is \cref{lem:gau-colrows}; the equation of $\Phi_0$ is that of \cref{lem:gau-HPF}~(i), (ii), which hold in the same way for coefficients in $\bbQ(\zeta)\fps q$; item (iii) there concerns coefficients in $\bbQ\fps u$ only.

\begin{thm}[Flatness of the $R$-matrix and its value at $q=0$, (QRR); quoted]\genusstatus{AI-proposed, unchecked}
\label{thm:set-qrr}
For a finite multiset $\mathbf w$ of non-zero numbers put
\begin{equation}
\label{eq:set-hw}
\hw(z)\coloneqq\exp\Bigl(-\sum_{k\geq1}\bk_k\,z^{2k-1}\sum_{a\in\mathbf w}a^{1-2k}\Bigr),\qquad\bk_k\coloneqq\frac{B_{2k}}{2k(2k-1)},
\end{equation}
and $\bw_i\coloneqq\{\zeta^i-\zeta^j\}_{j\neq i}\cup\{-5\zeta^i\}$. The $R$-matrix of \cref{def:set-rows} has coefficients in $\End(\bbH)\otimes\bbQ(\zeta)\fps q$ and satisfies \cref{def:set-rows}~(b); and $R(q=0)=\diag(h_{\bw_0},\dots,h_{\bw_4})$ in the basis $\tilde e_i(q=0)$: $R(0,z)\tilde e_i(0)=h_{\bw_i}(z)\,\tilde e_i(0)$, where $R(0,z)$ is the part of order $q^0$ of $R(z)$, to all orders in $z$.
\end{thm}

Flatness is \cite[Lemma~5.2, (16)]{GJR18}: there $R_{i\bar\alpha}=\eta(H^i,R\,\bar e_\alpha)$ is our row $\Qr^{(\alpha)}_i$, $L_\alpha=\zeta^\alpha L$ at $\lambda=1$, and the matrix on the right is that of $C_1H*$, so that (16) is literally \cref{prop:set-cyclicLR}~(iii); see also \cite[\S4.2, (12)--(13)]{GJR17}, and in general \cite[Section~1, Proposition (a)--(c)]{Giv01} and \cite[Props.~7.7, 8.2~(ii), 8.4, 8.7~(ii)]{MR2917177}. The value at $q=0$ is \cite[proof of Lemma~5.2, (17)]{GJR18} and \cite[proof of Lemma~4.6, (19)--(20)]{GJR17} at $t=0$; both refer to \cite[Section~2.2, Theorem~1 and the Remark after it]{Giv01} and \cite[Theorem~1, Corollary~1]{MR2276766}. The value is used only in the following form, for every colour $i$:
\begin{equation}
\label{eq:set-qrr}
\Phi_i(z,1)=h_{\bw_i}(z);
\end{equation}
in particular, with $h\coloneqq h_{\bw_0}$ and $h_n\coloneqq\coef{z^n}h$, $A_n(q=0)=\mathrm{ph}_n(1)=h_n$.

\begin{proof}[Proof of \eqref{eq:set-qrr}]
At $q=0$, $u=1$, $L=1$, $\zh=z$ and $C_0=1$. For every $q$, $\eta(1,\tilde e_i)=\Delta_i^{1/2}\eta(e_i,e_i)=\Delta_i^{-1/2}=\zeta^iL/C_0$ (\cref{cor:set-cyclicPF5}~(a), \cref{def:set-frob}). By (QRR) the part of order $q^0$ of $\Qr^{(i)}_0$ is $h_{\bw_i}(z)\,\eta(1,\tilde e_i)|_{q=0}=\zeta^ih_{\bw_i}(z)$, so $\Phi_i(z,1)=h_{\bw_i}(z)$.
\end{proof}

\begin{rmk}[The twist sign]
\label{rmk:set-qrrsign}
The weight $-5\zeta^i$ in $\bw_i$ is minus the fibre weight $5\zeta^i$ of $\cO(5)$ at $p_i$, the sign of an Euler-class twist: the Chern root $5\zeta^i$ contributes $\exp\bigl(+\sum_m\bk_m(z/5\zeta^i)^{2m-1}\bigr)=\exp\bigl(-\sum_m\bk_mz^{2m-1}(-5\zeta^i)^{1-2m}\bigr)$. The same sign appears in \cite[(17)]{GJR18}, \cite[Corollary~1]{MR2276766} and \cite[(26)]{LP20}. Pairing $j$ with $5-j$ gives $\frac1{1-\zeta^j}+\frac1{1-\zeta^{-j}}=1$, so $\sum_{a\in\bw_0}a^{-1}=2-\frac15=\frac95$ and $h_1=-\frac1{12}\cdot\frac95=-\frac3{20}$, whereas the weight $+5$ would give $-\frac{11}{60}$. The value $h_1=-\frac3{20}$ comes from (QRR). It agrees with \cref{lem:gau-HPF}~(iii), which forces $\mathrm{ph}_1=-\frac3{20}$ for a solution with coefficients in $\bbQ\fps u$; but this is not an independent check of the sign: over $\bbQ\fps q$ the equation $(4-5\theta_u)\mathrm{ph}_1=-\frac35$ of \cref{lem:gau-HPF}~(ii) also has the solutions $-\frac3{20}+k\,u^{4/5}$, $u^{4/5}=(1-5^5q)^{4/5}$ (the freedom of \cref{lem:gau-wkb}~(ii)), and that the actual $\mathrm{ph}_1$ is the constant one is part of Theorem~C, which uses (QRR).
\end{rmk}

\begin{lem}[Rotation of the colours]\genusstatus{AI-proposed, unchecked}
\label{lem:set-hrot}
$h_{\bw_i}(z)=h_{\bw_0}(\zeta^{-i}z)$ for every $i$.
\end{lem}

\begin{proof}
As multisets $\bw_i=\zeta^i\bw_0$, so $\sum_{a\in\bw_i}a^{1-2k}z^{2k-1}=\sum_{a\in\bw_0}a^{1-2k}(\zeta^{-i}z)^{2k-1}$ in \eqref{eq:set-hw}.
\end{proof}

\begin{lem}[Flatness and the value at $q=0$ determine $R$]\genusstatus{AI-proposed, unchecked}
\label{lem:set-Runique}
Assume (M) and (G0). There is exactly one series $R(z)=1+\sum_{k\geq1}R_kz^k$, $R_k\in\End(\bbH)\otimes\bbQ(\zeta)\fps q$, which satisfies \cref{def:set-rows}~(b) and $R(0,z)\tilde e_i(0)=h_{\bw_i}(z)\tilde e_i(0)$ for every $i$; and it satisfies \cref{def:set-rows}~(a).
\end{lem}

\begin{proof}
The $\tilde e_i$ have flat components in $\bbQ(\zeta)\fps q$ (\cref{cor:set-cyclicPF5}~(a), \cref{def:set-frob}: $\kappa^{\mathrm{LR}}_b$, $C_1/L$ and $C_0/L$ lie in $\bbQ\fps q$) and are orthonormal, so they form a basis of $\bbH\otimes\bbQ(\zeta)\fps q$ over $\bbQ(\zeta)\fps q$, and $x=\sum_j\eta(\tilde e_j,x)\tilde e_j$. Write $R(z)\tilde e_i=\sum_jY_{ji}(z)\tilde e_j$, $Y=1+\sum_{k\geq1}Y_kz^k$ with entries in $\bbQ(\zeta)\fps q$, and $\Dq\tilde e_i=\sum_j\Theta_{ji}\tilde e_j$, $\Theta_{ji}=\eta(\tilde e_j,\Dq\tilde e_i)$ ($\Theta$ a local letter; $\Dq$ acts on the flat components). As $\eta(\tilde e_i,\tilde e_j)=\delta_{ij}$ and the pairing of the flat basis are constant, $\Theta$ is antisymmetric, and its entries lie in $q\bbQ(\zeta)\fps q$, because $\Dq$ maps $\bbQ(\zeta)\fps q$ onto $q\bbQ(\zeta)\fps q$, with kernel $\bbQ(\zeta)$. Since $C_1H*\tilde e_j=\zeta^jL\,\tilde e_j$ (\cref{cor:set-cyclicPF5}~(a)), property (b), multiplied by $C_1$, reads $z\Dq Y+z\Theta Y=[\Dq U,Y]$ with $\Dq U\coloneqq\diag(\Dq u_j)=\diag(\zeta^jL)$, that is,
\[
(\Dq Y_k+\Theta Y_k)_{ji}=(\zeta^j-\zeta^i)L\,(Y_{k+1})_{ji}\qquad(k\geq0,\ Y_0=1).
\]
(The coefficient of $z^0$ is $[\Dq U,Y_0]=0$.) For $j\neq i$ this determines the off-diagonal entries of $Y_{k+1}$ from $Y_k$, as $(\zeta^j-\zeta^i)L$ is a unit of $\bbQ(\zeta)\fps q$ ($L\in1+q\bbQ\fps q$); they lie in $q\bbQ(\zeta)\fps q$, because the left side does. For $j=i$, with $k+1$ in place of $k$, it reads $\Dq(Y_{k+1})_{ii}=-\sum_{l\neq i}\Theta_{il}(Y_{k+1})_{li}$. The right side lies in $q\bbQ(\zeta)\fps q$, so this equation has a solution in $\bbQ(\zeta)\fps q$ for every prescribed constant term, and exactly one; the equation with $j=i$ and $k=0$ holds, as $\Theta_{ii}=0$. So, by induction on $k$, for every family of constants $c_{k,i}\in\bbQ(\zeta)$ ($k\geq1$) there is exactly one solution $Y$ with $(Y_k)_{ii}(q=0)=c_{k,i}$, and every solution has $Y(q=0)$ diagonal. The condition at $q=0$ is $Y(q=0)=\diag(h_{\bw_i})$, that is $c_{k,i}=\coef{z^k}h_{\bw_i}$; this proves existence and uniqueness, $R(z)$ being the endomorphism with the matrix $Y(z)$ in the basis $\tilde e_i$. For (a) put $N(z)\coloneqq Y(-z)^TY(z)$ (a local letter). From the equation and its transpose at $-z$, using $\Theta^T=-\Theta$, $z\Dq N=[\Dq U,N]$. For $j\neq i$ the lowest coefficient in $z$ of $N_{ji}$ is annihilated by the unit $(\zeta^j-\zeta^i)L$, so $N_{ji}=0$; and $\Dq N_{ii}=0$, so $N_{ii}\in\bbQ(\zeta)\fps z$ is its value at $q=0$, which is $h_{\bw_i}(-z)h_{\bw_i}(z)=1$ because $Y(q=0)$ is diagonal and the exponent of $h_{\bw_i}$ is odd in $z$ \eqref{eq:set-hw}. Hence $Y(-z)^TY(z)=1$, which is (a) in the orthonormal basis $\tilde e_i$.
\end{proof}

%-------------------------------------------------------------------------------
\subsection{The graph sum}
\label{sec:set-graph}

\begin{defn}[The Givental--Teleman graph sum]
\label{def:set-graphsum}
Let $\Omega$ be a cohomological field theory with unit $1$ \cite[Def.~0.3]{PPZ15} on a free module $V$ of finite rank with a perfect symmetric pairing $\eta$ (its Gram matrix is invertible) over a commutative $\bbQ$-algebra, and $R(z)=1+O(z)\in\End(V)\fps z$ symplectic, $R(-z)^*R(z)=1$ with ${}^*$ the $\eta$-adjoint \cite[(5)]{PPZ15}. The ancestor correlators are $\langle\psi^{k_1}a_1\cdots\psi^{k_n}a_n\rangle_h\coloneqq\int_{\overline{\mathcal M}_{h,n}}\Omega_{h,n}(a_1,\dots,a_n)\prod_j\psi_j^{k_j}$; for $q\in V\fps z$ (a point of the loop space, not the parameter $q$) the dilaton-shifted potential is $\overline F^\Omega_h(q)\coloneqq\sum_{n:\,2h-2+n>0}\frac1{n!}\langle(q+z1)^n\rangle_h$, with $z^ka\mapsto\psi^ka$ (base point $q=-z1$), and $\dd^n\overline F^\Omega_h(q)$ is the symmetric form $(x_1,\dots,x_n)\mapsto\sum_m\frac1{m!}\langle x_1\cdots x_n(q+z1)^m\rangle_h$ on $V\fps z$. The \emph{edge} and the \emph{leaf point} of $R$ are
\begin{gather*}
E_R(\psi,\psi')\coloneqq\frac{\eta^{-1}-\bigl(R^{-1}(\psi)\otimes R^{-1}(\psi')\bigr)\eta^{-1}}{\psi+\psi'}\in(V\otimes V)\fps{\psi,\psi'},\\
q_R(z)\coloneqq R^{-1}(z)(-z1)=-z1+\ttr_R(z),
\end{gather*}
where $\eta^{-1}\in V\otimes V$ is the inverse bivector, $E_R$ is a power series by the symplectic condition \cite[Rem.~2.1, (6)]{PPZ15}, and $\ttr_R(z)\coloneqq z1-zR^{-1}(z)1\in z^2V\fps z$ is the translation of \cite[Def.~2.13]{PPZ15}. For a stable graph $\Gamma$ without legs, with vertices $v$ of genus $g_v$ and valence $n_v$ and of total genus $g=\sum_vg_v+h^1(\Gamma)$, and a bivector series $E$ symmetric under the exchange of its two slots (together with $\psi$, $\psi'$), the contraction
\[
\mathrm{Cont}_\Gamma[E,q;\Omega]\coloneqq|\Aut\Gamma|^{-1}\Bigl\langle\prod_v\dd^{n_v}\overline F^\Omega_{g_v}(q),\ \prod_eE\Bigr\rangle
\]
contracts, at each vertex, the $n_v$-th derivative of $\overline F^\Omega_{g_v}$ at $q$, one slot for each half-edge, with one copy of $E(\psi,\psi')$ for each edge, its two slots at the two half-edges; sums over graphs carry no further factor. For $q=q_R$ the contraction is finite, as $\ttr_R\in z^2V\fps z$. Let $R.\Omega$ be the unit-preserving action of \cite[Defs.~2.2, 2.5, 2.13]{PPZ15}: legs $R^{-1}(\psi)$, edges $E_R$, translation $\ttr_R$, as in Givental and Teleman \cite[p.~20]{PPZ15}. There the translation points are pushed forward, whereas the vertex tensors above insert them; for $g\geq2$ the two forms agree (\cref{lem:set-pushins}~(iii)). The sum $\sum_\Gamma\mathrm{Cont}_\Gamma[E_R,q_R;\Omega]$ over the leg-free stable graphs of genus $g$ is the \emph{graph sum of genus $g$}. The definitions of \cite{PPZ15} used in the note are finite sums; they make sense over every commutative $\bbQ$-algebra $A$, for $V$ free of finite rank with a perfect pairing and classes in $H^*(\overline{\mathcal M}_{g,n},\bbQ)\otimes_\bbQ A$. Two propositions of \cite{PPZ15}, stated there over $\bbQ$, are used over other rings: \cite[Prop.~2.12]{PPZ15} ($R.\Omega$ is a cohomological field theory with the unit of $\Omega$) and \cite[Prop.~2.14]{PPZ15} (the unit-preserving action is a left group action). They are applied only to a topological field theory $\TFT_y$ on $A^r$ with an orthonormal basis $e_i$, $(\TFT_y)_{g,n}(e_i,\dots,e_i)=y_i^{-(2g-2+n)}$, zero on mixed colours, with the unit $\sum_iy_ie_i$, $y_i\in A^\times$, and to series $R=1+O(z)$ with $R(-z)^TR(z)=1$, that is, $R=\exp(r)$ with $r=\sum_{k\geq1}r_kz^k$ and $r_k^T=(-1)^{k+1}r_k$. In a $\bbQ$-basis of $H^*(\overline{\mathcal M}_{g,n},\bbQ)$, for fixed $(g,n)$, each conclusion used is a finite family of identities between Laurent polynomials over $\bbQ$ in the $y_i$ and in finitely many free coordinates of the $r_k$. \cite{PPZ15} gives them at all points with $y\in(\bbQ^\times)^r$ and rational $r_k$; as $\bbQ$ is infinite they are identities of Laurent polynomials, and so they hold over $A$.

\emph{Colours.} When $\Omega=\TFT$ (\cref{def:set-frob}), which vanishes on mixed colours, a \emph{coloured graph} is a stable graph with a colour at each vertex, and the term of $\Gamma$ is the sum over its colourings. For a constant isometry $A\colon(W,\eta_W)\to(V,\eta)$ and $A_*\Omega\coloneqq\Omega\circ(A^{-1})^{\otimes n}$, $A_*(R.\Omega)=(ARA^{-1}).(A_*\Omega)$, immediate from \cite[Defs.~2.2, 2.13]{PPZ15}; it gives the graph sum in the colour basis (\cref{sec:sgn}).
\end{defn}

\input{sections/s02_setup_5}



